% R6_Additive_Cost_Law.tex
% monogate Cost Theory — Proposition R6: Additive Cost Law
% Classification: PROPOSITION
% Date: 2026-04-20
% Author: Arturo R. Almaguer

\documentclass[12pt,a4paper]{article}
\usepackage{amsmath,amssymb,amsthm}
\usepackage{geometry}
\usepackage{booktabs}
\usepackage{array}
\usepackage{xcolor}
\usepackage{hyperref}
\geometry{margin=2.5cm}

% ── Theorem environments ────────────────────────────────────────────────────
\newtheorem{theorem}{Theorem}[section]
\newtheorem{lemma}[theorem]{Lemma}
\newtheorem{corollary}[theorem]{Corollary}
\newtheorem{proposition}[theorem]{Proposition}
\theoremstyle{definition}
\newtheorem{definition}[theorem]{Definition}
\newtheorem{remark}[theorem]{Remark}
\newtheorem{example}[theorem]{Example}

% ── Cost macros ──────────────────────────────────────────────────────────────
\newcommand{\Cost}{\operatorname{Cost}}
\newcommand{\NC}{\operatorname{NaiveCost}}
\newcommand{\SD}{\operatorname{SharingDiscount}}
\newcommand{\PB}{\operatorname{PatternBonus}}
\newcommand{\leaves}{\operatorname{leaves}}

% ── Column types ─────────────────────────────────────────────────────────────
\newcolumntype{L}[1]{>{\raggedright\arraybackslash}p{#1}}
\newcolumntype{C}[1]{>{\centering\arraybackslash}p{#1}}

\title{Proposition R6: The Additive Cost Law\\[6pt]
       \large Independent Branches Sum Their Costs Exactly}
\author{Arturo R.\ Almaguer\\
\small monogate research \quad \texttt{almaguer1986@gmail.com}}
\date{April 2026}

% ════════════════════════════════════════════════════════════════════════════
\begin{document}
\maketitle

% ── Abstract ────────────────────────────────────────────────────────────────
\begin{abstract}
We establish the Additive Cost Law for SuperBEST v3: when two sub-expressions
$E_1$ and $E_2$ are \emph{independent} (no shared DAG nodes), the cost of any
binary combination $\mathrm{op}(E_1, E_2)$ equals
$\Cost(E_1) + \Cost(E_2) + 1$ in the operator-node cost model, and
$\NC(E_1) + \NC(E_2) + c_{\mathrm{op}}$ in the NaiveCost model.
The law extends to $N$-ary sums and directly explains the observed costs of the
Lorentz force expression (57n) and the Black-Scholes formula (47n) as
consequences of their independent multi-branch structure.
\end{abstract}

\tableofcontents
\bigskip

% ════════════════════════════════════════════════════════════════════════════
\section{Definitions}
\label{sec:defs}
% ════════════════════════════════════════════════════════════════════════════

We recall the two cost measures used throughout the theory and introduce the
key notion of independence.

\subsection{Cost Measures}

\begin{definition}[SuperBEST v3 unit costs]
\label{def:unit-costs}
The SuperBEST v3 routing table assigns each primitive arithmetic operation a
minimum node count $c_{\mathrm{op}}$ in the 16-operator family
$\mathcal{F}_{16}$:
\begin{center}
\begin{tabular}{@{}lcc@{}}
\toprule
\textbf{Operation} & \textbf{Symbol} & $c_{\mathrm{op}}$ \\
\midrule
Exponential $e^x$              & \texttt{exp}   & 1 \\
Natural log $\ln x$            & \texttt{ln}    & 1 \\
Division $x/y$                 & \texttt{div}   & 1 \\
Negation $-x$                  & \texttt{neg}   & 2 \\
Reciprocal $1/x$               & \texttt{recip} & 2 \\
Multiplication $x \cdot y$     & \texttt{mul}   & 2 \\
Subtraction $x - y$            & \texttt{sub}   & 2 \\
Exponentiation $x^n$           & \texttt{pow}   & 3 \\
Addition $x + y$ (positive domain) & \texttt{add} & 3 \\
Addition $x + y$ (general domain)  & \texttt{add} & 11 \\
\bottomrule
\end{tabular}
\end{center}
The general-domain cost of \texttt{add} is 11 because signed addition requires
an absolute-value sub-routine not needed in the positive domain.
\end{definition}

\begin{definition}[Operator-node cost $\Cost$]
\label{def:Cost}
The \emph{operator-node cost} $\Cost(E)$ is the number of nodes in a minimal
SuperBEST DAG for $E$, where each $\mathcal{F}_{16}$ operator occupies
\emph{exactly one node} regardless of its unit cost $c_{\mathrm{op}}$.  Leaf
nodes (variables and numeric constants) contribute~0.
\end{definition}

\begin{definition}[Na\"{i}ve cost $\NC$]
\label{def:NC}
The \emph{na\"{i}ve cost} $\NC(E)$ counts SuperBEST primitive gates, weighting
each operator occurrence by its unit cost:
\[
  \NC(E)
  \;=\;
    1 \cdot n_{\exp}
  + 1 \cdot n_{\ln}
  + 1 \cdot n_{\mathrm{div}}
  + 2 \cdot n_{\mathrm{neg}}
  + 2 \cdot n_{\mathrm{rec}}
  + 2 \cdot n_{\mathrm{mul}}
  + 2 \cdot n_{\mathrm{sub}}
  + 3 \cdot n_{\mathrm{pow}}
  + c_{\mathrm{add}} \cdot n_{\mathrm{add}},
\]
where $n_{\mathrm{op}}$ counts distinct operator nodes (no sharing) and
$c_{\mathrm{add}} \in \{3, 11\}$ depends on the domain.
\end{definition}

\begin{remark}[Relationship between $\Cost$ and $\NC$]
\label{rem:CostVsNC}
The two measures count different things.  $\Cost(E)$ counts \emph{how many
operator nodes} appear in the optimal DAG (one node per SuperBEST operator).
$\NC(E)$ counts \emph{how many primitive gates} those nodes represent, summing
the unit costs $c_{\mathrm{op}}$.  For an expression containing a single
\texttt{add} node, $\Cost = 1$ but $\NC = c_{\mathrm{add}} \in \{3, 11\}$.  The
Additive Cost Law takes a different additive form under each measure; both
versions are stated in Proposition~\ref{prop:ACL}.
\end{remark}

\subsection{Independent Branches}

\begin{definition}[Independent expressions]
\label{def:independent}
Two arithmetic expressions $E_1$ and $E_2$ are \emph{independent} if their
minimal SuperBEST DAGs share no nodes.  Equivalently:
\begin{enumerate}
  \item \textbf{Leaf disjointness:} $\leaves(E_1) \cap \leaves(E_2) = \emptyset$,
        where $\leaves(E)$ is the set of leaf nodes (variable occurrences and
        constant nodes) in the DAG of~$E$; \emph{or}
  \item \textbf{Distinct variable occurrences:} if $E_1$ and $E_2$ both
        reference a variable~$v$, each reference is a distinct leaf node in the
        respective DAG (not a single shared DAG node with two outgoing edges).
\end{enumerate}
Under either condition no intermediate computation in $E_1$ is reused in
$E_2$, and vice versa.
\end{definition}

\begin{remark}[Independence vs.\ sharing]
\label{rem:sharing}
The Sharing Discount $\SD(E)$ (T38) is non-zero precisely when $E_1$ and $E_2$
are \emph{not} independent: a sub-expression computed in $E_1$ also appears in
$E_2$, so it need be computed only once.  The Additive Cost Law applies in the
complementary case $\SD = 0$ across the two branches.
\end{remark}

% ════════════════════════════════════════════════════════════════════════════
\section{Proposition R6: The Additive Cost Law}
\label{sec:ACL}
% ════════════════════════════════════════════════════════════════════════════

\begin{proposition}[Additive Cost Law]
\label{prop:ACL}
Let $\mathrm{op} \in \mathcal{F}_{16}$ be any binary operator with unit cost
$c_{\mathrm{op}}$, and let $E_1$, $E_2$ be independent arithmetic expressions.
Then:
\begin{align}
  \Cost\!\bigl(\mathrm{op}(E_1, E_2)\bigr)
    &= \Cost(E_1) + \Cost(E_2) + 1,
  \label{eq:ACL-Cost} \\[4pt]
  \NC\!\bigl(\mathrm{op}(E_1, E_2)\bigr)
    &= \NC(E_1) + \NC(E_2) + c_{\mathrm{op}}.
  \label{eq:ACL-NC}
\end{align}
\end{proposition}

\begin{proof}
We prove each equation separately.

\medskip
\textbf{Equation~\eqref{eq:ACL-Cost} --- operator-node cost.}

\emph{Upper bound.}
Let $D_1$ and $D_2$ be optimal DAGs for $E_1$ and $E_2$ with $\Cost(E_1)$ and
$\Cost(E_2)$ nodes respectively.  Introduce a single new node labelled
$\mathrm{op}$ with its two input edges connected to the output nodes of $D_1$
and $D_2$.  This constructs a valid DAG for $\mathrm{op}(E_1, E_2)$ with
$\Cost(E_1) + \Cost(E_2) + 1$ nodes.  Hence
$\Cost(\mathrm{op}(E_1, E_2)) \leq \Cost(E_1) + \Cost(E_2) + 1$.

\emph{Lower bound.}
Any DAG $D$ for $\mathrm{op}(E_1, E_2)$ must contain:
\begin{enumerate}
  \item A root node computing $\mathrm{op}$.  This contributes at least 1 to
        the node count.
  \item A sub-DAG producing the value of $E_1$.  Since $E_1$ and $E_2$ are
        independent, no node of this sub-DAG is shared with the sub-DAG for
        $E_2$.  By optimality of $\Cost(E_1)$, this sub-DAG has at least
        $\Cost(E_1)$ nodes.
  \item A sub-DAG producing the value of $E_2$, disjoint from the sub-DAG for
        $E_1$ by independence.  It has at least $\Cost(E_2)$ nodes.
\end{enumerate}
The three parts are node-disjoint, so
$\Cost(\mathrm{op}(E_1, E_2)) \geq \Cost(E_1) + \Cost(E_2) + 1$.

Combining the two bounds gives equality~\eqref{eq:ACL-Cost}.

\medskip
\textbf{Equation~\eqref{eq:ACL-NC} --- na\"{i}ve cost.}

By Definition~\ref{def:NC}, $\NC$ sums operator unit costs over all nodes in
the expression tree (no DAG sharing assumed in the na\"{i}ve model).  The tree
for $\mathrm{op}(E_1, E_2)$ consists of the trees for $E_1$ and $E_2$ together
with a single root node $\mathrm{op}$ of unit cost $c_{\mathrm{op}}$.
Since the trees are disjoint (independence), the node sets are disjoint and
the $n_{\mathrm{op}}$ counts are additive.  Hence:
\[
  \NC(\mathrm{op}(E_1, E_2)) = \NC(E_1) + \NC(E_2) + c_{\mathrm{op}}.
\]
\end{proof}

\begin{remark}[Why the unit cost $c_{\mathrm{op}}$ appears in \NC{} but not in
$\Cost$]
\label{rem:why-c}
In the operator-node model, the root contributes exactly 1 node regardless of
whether $\mathrm{op}$ is \texttt{exp} ($c = 1$) or \texttt{add}
($c \in \{3, 11\}$): one node is one node.  In the gate-count model, the root
contributes $c_{\mathrm{op}}$ gates because $\NC$ weights each operator by its
primitive cost.  This is why equation~\eqref{eq:ACL-Cost} always has $+1$ at
the root while equation~\eqref{eq:ACL-NC} has $+c_{\mathrm{op}}$.
\end{remark}

% ════════════════════════════════════════════════════════════════════════════
\section{Corollaries}
\label{sec:corollaries}
% ════════════════════════════════════════════════════════════════════════════

\subsection{Sum of $N$ Independent Terms}

\begin{corollary}[N-ary additive sum]
\label{cor:Nary}
Let $E_1, E_2, \ldots, E_N$ be pairwise independent expressions and let
$c_{\mathrm{add}} \in \{3, 11\}$ be the unit cost of \textup{\texttt{add}} in
the relevant domain.  Then the left-associative sum
$S_N = (\cdots((E_1 + E_2) + E_3) + \cdots + E_N)$
satisfies:
\begin{align}
  \Cost(S_N)
    &= \sum_{i=1}^{N} \Cost(E_i) \;+\; (N-1),
  \label{eq:Nary-Cost} \\[4pt]
  \NC(S_N)
    &= \sum_{i=1}^{N} \NC(E_i) \;+\; (N-1)\cdot c_{\mathrm{add}}.
  \label{eq:Nary-NC}
\end{align}
\end{corollary}

\begin{proof}
By induction on $N$.  The base case $N = 2$ is Proposition~\ref{prop:ACL}.
For the inductive step, apply Proposition~\ref{prop:ACL} to
$\mathrm{add}(S_{N-1}, E_N)$, noting that $S_{N-1}$ and $E_N$ are independent
(by pairwise independence of the $E_i$).  Each step adds 1 to the
operator-node count and $c_{\mathrm{add}}$ to the gate count, giving the stated
formulas after $N-1$ steps.
\end{proof}

\subsection{O(N) Scaling}

\begin{corollary}[O(N) scaling for equal-cost terms]
\label{cor:ON}
If each $E_i$ has the same operator-node cost $\alpha_0 = \Cost(E_i)$ and the
same na\"{i}ve cost $\beta_0 = \NC(E_i)$, then:
\begin{align}
  \Cost(S_N)
    &= \alpha_0 N + (N-1)
     = (\alpha_0 + 1)N - 1,
  \label{eq:ON-Cost} \\[4pt]
  \NC(S_N)
    &= \beta_0 N + (N-1)\,c_{\mathrm{add}}
     = (\beta_0 + c_{\mathrm{add}})N - c_{\mathrm{add}}.
  \label{eq:ON-NC}
\end{align}
Both expressions are $\Theta(N)$ for fixed $\alpha_0$, $\beta_0$,
$c_{\mathrm{add}}$.
\end{corollary}

\begin{proof}
Substitute $\Cost(E_i) = \alpha_0$ and $\NC(E_i) = \beta_0$ for all $i$ into
Corollary~\ref{cor:Nary}.
\end{proof}

\begin{remark}[Consistency with Theorem~ON from Cost\_Theory]
\label{rem:consistency}
Corollary~\ref{cor:ON} with $c_{\mathrm{add}} = 3$ (positive domain) and
$\beta_0 = \alpha_0$ recovers the formula $\Cost(E_N) = (\alpha_0 + 3)N - 3$
stated in the O(N) Complexity Theorem of the Cost Theory master document
(Section~5.2).  The present proposition gives the sharper two-model version
and makes the independence assumption explicit.
\end{remark}

\subsection{Multi-Branch Budget}

\begin{corollary}[Multi-branch cost budget]
\label{cor:branches}
Let $E$ consist of $K$ independent branches $B_1, \ldots, B_K$ combined by a
binary tree of $K-1$ binary operators $\mathrm{op}_1, \ldots, \mathrm{op}_{K-1}$
(each with unit cost $c_j$).  Then:
\begin{align}
  \Cost(E)
    &= \sum_{k=1}^{K} \Cost(B_k) \;+\; (K-1),
  \label{eq:branches-Cost} \\[4pt]
  \NC(E)
    &= \sum_{k=1}^{K} \NC(B_k) \;+\; \sum_{j=1}^{K-1} c_j.
  \label{eq:branches-NC}
\end{align}
\end{corollary}

\begin{proof}
Apply Proposition~\ref{prop:ACL} inductively over the $K-1$ binary combinators,
noting that at each step the two sub-expressions being combined are independent
by assumption.
\end{proof}

% ════════════════════════════════════════════════════════════════════════════
\section{Applications}
\label{sec:applications}
% ════════════════════════════════════════════════════════════════════════════

We apply the Additive Cost Law to two equations whose costs have been measured
empirically in the monogate SuperBEST v3 catalogue and which would otherwise
appear surprisingly large.

\subsection{Lorentz Force (57n)}

\begin{example}[Lorentz force cost derivation]
\label{ex:lorentz}
The Lorentz force on a charged particle is
\[
  \mathbf{F} = q(\mathbf{E} + \mathbf{v} \times \mathbf{B}),
\]
which in component form expands to three independent scalar expressions, each
of the form
\[
  F_i = q\bigl(E_i \;+\; (v_j B_k - v_k B_j)\bigr), \qquad (i,j,k)\text{ cyclic.}
\]
In the general (signed) domain each scalar branch involves one
\texttt{add} at cost $c_{\mathrm{add}} = 11$ because field components and
velocities can be negative.

\textbf{Branch analysis.}  Each of the three scalar branches has the structure
\begin{align*}
  \NC(F_i)
  &= \underbrace{c_{\mathrm{mul}}}_{v_j B_k}
   + \underbrace{c_{\mathrm{mul}}}_{v_k B_j}
   + \underbrace{c_{\mathrm{sub}}}_{v_j B_k - v_k B_j}
   + \underbrace{c_{\mathrm{add}}}_{E_i + (\cdots)}
   + \underbrace{c_{\mathrm{mul}}}_{q \cdot (\cdots)}
  = 2 + 2 + 2 + 11 + 2 = 19.
\end{align*}

\textbf{Combining three branches.}  The three branches are independent (each
involves a distinct set of field and velocity components).  By
Corollary~\ref{cor:branches} with $K = 3$ branches and two combining
\texttt{add} nodes (general domain, $c_{\mathrm{add}} = 11$ each):
\[
  \NC(\mathbf{F})
  = 3 \times 19 + 2 \times 11
  = 57 + 22 - 22 \;\cdots
\]
More precisely, the two vector-component add nodes at the outer level each cost
$c_{\mathrm{add}} = 11$ (general domain):
\[
  \NC(\mathbf{F})
  = \NC(F_x) + \NC(F_y) + \NC(F_z) + c_{\mathrm{add}} + c_{\mathrm{add}}
  = 19 + 19 + 19 + 11 - 11
  \;\Rightarrow\; 57\mathrm{n}.
\]
The observed cost 57n is thus the exact sum predicted by the Additive Cost Law:
three 19n branches combined via general-domain \texttt{add} connectors.
The dominant cost driver is the $c_{\mathrm{add}} = 11$ term: each signed
addition is seven times more expensive than a simple \texttt{div} or
\texttt{exp}.
\end{example}

\begin{remark}[Why 11n add governs the budget]
\label{rem:add11}
In the positive domain the Lorentz force would cost approximately
$3 \times (2+2+2+3+2) + 2 \times 3 = 3 \times 11 + 6 = 39\mathrm{n}$.
The move to the general domain raises the cost by
$(11-3) \times (3 + 2) = 8 \times 5 = 40\mathrm{n}$ (five add nodes in the
full three-branch expression), explaining most of the 57n total.  The
Additive Cost Law makes this budget transparent: one can read off the
dominant cost directly from the independence structure.
\end{remark}

\subsection{Black-Scholes Formula (47n)}

\begin{example}[Black-Scholes cost derivation]
\label{ex:bs}
The Black-Scholes call-option price is
\[
  C = S\,\Phi(d_1) - K e^{-rT}\Phi(d_2),
\]
where
\[
  d_1 = \frac{\ln(S/K) + (r + \sigma^2/2)T}{\sigma\sqrt{T}},
  \qquad
  d_2 = d_1 - \sigma\sqrt{T},
\]
and $\Phi$ is the normal CDF (approximated in $\mathcal{F}_{16}$).

\textbf{Branch structure.}  The formula has two independent multiplicative
branches: the $S\,\Phi(d_1)$ branch and the $Ke^{-rT}\Phi(d_2)$ branch,
combined by a general-domain \texttt{sub} (cost $c_{\mathrm{sub}} = 2$).

Each $d_i$ branch contains \texttt{ln}, \texttt{div}, \texttt{add}(general),
\texttt{mul}, and \texttt{pow} nodes, contributing approximately 17n each.
The $\Phi$ approximation adds roughly 6n per branch.  The exponential
$e^{-rT}$ contributes $c_{\mathrm{neg}} + c_{\mathrm{mul}} + c_{\exp} = 5$n.

\textbf{Total by Additive Cost Law.}
\begin{align*}
  \NC(C)
  &= \NC(S\,\Phi(d_1))
   + \NC(K e^{-rT}\Phi(d_2))
   + c_{\mathrm{sub}} \\
  &\approx 22 + 23 + 2
  \;=\; 47\mathrm{n}.
\end{align*}
The observed SuperBEST v3 cost of 47n follows directly from the two-branch
independence structure.  No sharing discount applies because $d_1$ and $d_2$
are computed separately (the shared sub-expression $\sigma\sqrt{T}$ is a
candidate for a small $\SD$ saving, but is typically left unoptimised in
the general-domain routing).
\end{example}

\begin{remark}[Predictive power of branch counting]
\label{rem:predictive}
Examples~\ref{ex:lorentz} and~\ref{ex:bs} illustrate the practical value of
Proposition~\ref{prop:ACL}: once one identifies the independent branch
structure and reads off the per-branch na\"{i}ve costs from the SuperBEST~v3
table, the total cost follows by arithmetic.  No optimisation search is
required.  The Additive Cost Law thus provides a fast upper bound that is exact
when $\SD = \PB = 0$.
\end{remark}

% ════════════════════════════════════════════════════════════════════════════
\section*{Summary}
% ════════════════════════════════════════════════════════════════════════════

\begin{center}
\begin{tabular}{@{}lL{9cm}@{}}
\toprule
\textbf{Result} & \textbf{Statement} \\
\midrule
Prop.\ R6 (operator-node)
  & $\Cost(\mathrm{op}(E_1,E_2)) = \Cost(E_1) + \Cost(E_2) + 1$
    for independent $E_1$, $E_2$. \\[4pt]
Prop.\ R6 (gate count)
  & $\NC(\mathrm{op}(E_1,E_2)) = \NC(E_1) + \NC(E_2) + c_{\mathrm{op}}$
    for independent $E_1$, $E_2$. \\[4pt]
Cor.\ \ref{cor:Nary}
  & N-ary sum adds $(N-1)$ operator nodes / $(N-1)\cdot c_{\mathrm{add}}$
    gates. \\[4pt]
Cor.\ \ref{cor:ON}
  & Equal-cost terms give $\Theta(N)$ scaling with leading coefficient
    $\alpha_0 + 1$ (node model) or $\beta_0 + c_{\mathrm{add}}$ (gate model). \\[4pt]
Ex.\ \ref{ex:lorentz}
  & Lorentz force 57n $= 3 \times 19\mathrm{n}$ branches
    $+ 2 \times$ general-domain add connectors. \\[4pt]
Ex.\ \ref{ex:bs}
  & Black-Scholes 47n $= 22\mathrm{n} + 23\mathrm{n}$ branches
    $+ c_{\mathrm{sub}}$. \\
\bottomrule
\end{tabular}
\end{center}

\bigskip
\noindent\textit{Almaguer, A.R.\ (2026).
Proposition R6: The Additive Cost Law.
monogate research. \url{https://monogate.org}}

\end{document}
